% Corte mecánico íntegro: parte_iv.tex:1--54.
% Residencia material del ensamblaje sucesor de 102 capítulos.
% Residencia de realización: capítulo 67.
\part{Géométrie différentielle et réalisations dimensionnelles}

\chapter{Courbure et torsion des courbes publiées}

\section{Courbe plane polaire}

Soit
\[
 \gamma(\theta)=r(\theta)(\cos\theta,\sin\theta).
\]
Les primes désignant les dérivées par rapport à \(\theta\), la vitesse et la courbure sont
\begin{align}
 v(\theta)&=\sqrt{r^2+(r')^2},\\
 \kappa_{\mathrm{F}}(\theta)&=
 \frac{r^2+2(r')^2-rr''}
 {\bigl(r^2+(r')^2\bigr)^{3/2}}.
 \label{espiral:eq:curvatura-polar}
\end{align}
La longueur d'arc entre \(\theta_0\) et \(\theta_1\) est
\[
 s=\int_{\theta_0}^{\theta_1}
 \sqrt{r^2+(r')^2}\,\dd\theta.
\]
Ces grandeurs appartiennent à la courbe publiée. Elles ne sont ni le régime TRIT ni le cocycle radial.

\section{Spécialisations}

Pour la spirale logarithmique \(r=r_0\e^{b\theta}\),
\[
 r'=br,\qquad r''=b^2r,
\]
et
\begin{equation}
 \kappa_{\mathrm{F}}(\theta)
 =\frac{1}{r(\theta)\sqrt{1+b^2}}.
 \label{espiral:eq:curvatura-log}
\end{equation}
L'angle \(\psi\) entre le vecteur tangent et la direction radiale satisfait
\[
 \tan\psi=\frac{r}{r'}=\frac1b,
\]
et est donc constant. C'est la propriété équiangulaire de la spirale logarithmique.


Pour la spirale d'Archimède \(r=a\theta\),
\[
 \kappa_{\mathrm{F}}(\theta)
 =\frac{\theta^2+2}{a(\theta^2+1)^{3/2}}.
\]
Pour Fermat \(r=a\sqrt\theta\), la singularité de l'origine exige de déclarer \(\theta>0\). Pour les branches hyperbolique et lituus, l'origine angulaire est une asymptote et non un point du domaine.
